\originalchapter{自旋与二态系统}
\originalsection{自旋二分之一}
Stern--Gerlach装置让原子束穿过不均匀磁场, 再在屏上记录偏转. 磁矩的势能$H_{\rm mag}=-\boldsymbol\mu\cdot\boldsymbol B$, 沿$z$的场梯度产生$F_z\simeq\mu_z\partial_zB_z$, 把磁矩分量转换为位置带. 均匀场主要使方向进动, 本身不提供这个分束力. 用自旋$1/2$的有效自由度描述双输出, 磁矩与自旋成正比. Stern--Gerlach实验中的分束表明某些内禀角动量的分量取离散值. 自旋不是一个有确定表面速度的经典小球自转模型; 它是独立于轨道坐标的量子自由度. 对自旋$1/2$, 选$S_z$本征基$\ket{+},\ket{-}$, 本征值为$\pm\hbar/2$. 任意归一纯态是$\alpha\ket{+}+\beta\ket{-}$, $|\alpha|^2+|\beta|^2=1$.

先选$z+$, 再沿$x$测量得到两种结果各半; 筛选$x+$后再测$z$仍各半. 因此$x+$不是原$z+$标签, 可取$\ket{x;\pm}=(\ket+\pm\ket-)/\sqrt2$. 从$S_x=(\hbar/2)(\ket{x;+}\bra{x;+}-\ket{x;-}\bra{x;-})$恢复非对角矩阵, 它是测量基变化而非多加一个经典随机标签. $y$向上态为$(\ket++\ii\ket-)/\sqrt2$, 相对$\ii$相位使它与$x$测量不同. 三种自旋分量可表示为$S_j=(\hbar/2)\sigma_j$, 其中
\[
\sigma_x=\begin{pmatrix}0&1\\1&0\end{pmatrix},\quad
\sigma_y=\begin{pmatrix}0&-\ii\\\ii&0\end{pmatrix},\quad
\sigma_z=\begin{pmatrix}1&0\\0&-1\end{pmatrix}.
\]
令$\epsilon_{ijk}$为Levi--Civita符号: $\epsilon_{123}=1$, 交换任意两指标变号, 重复指标时为零. 直接矩阵乘法给$\sigma_i\sigma_j=\delta_{ij}I+\ii\sum_k\epsilon_{ijk}\sigma_k$, 所以$[S_i,S_j]=\ii\hbar\sum_k\epsilon_{ijk}S_k$. 对$S_z$确定的态, $S_x,S_y$一般不确定. 例如$\ket+$中$\langle S_x\rangle=0$, 但$\langle S_x^2\rangle=\hbar^2/4$, 所以$\Delta S_x=\hbar/2$, 不是零.

对单位向量$\boldsymbol n=(\sin\theta\cos\phi,\sin\theta\sin\phi,\cos\theta)$, $\theta$是相对$z$轴的极角,$\phi$是$xy$平面方位角; 纯态的方向指均值向量方向, 三分量并非都确定. 沿该方向的自旋是$S_{\boldsymbol n}=\boldsymbol n\cdot\boldsymbol S$. $+\hbar/2$本征态可取
$\ket{\boldsymbol n,+}=\cos(\theta/2)\ket+ +\e^{\ii\phi}\sin(\theta/2)\ket-$. 代入$\boldsymbol n\cdot\boldsymbol\sigma$验证其本征值为1. 例如$\theta=\pi/2,\phi=0$恢复$x+$, $\phi=\pi/2$恢复$y+$, 直接矩阵乘法给各对应本征值$+\hbar/2$. 从$z$向上态测量沿$\boldsymbol n$自旋, 向上概率为$\cos^2(\theta/2)$.

\originalsection{一般二态Hamilton算符}
二态不是“只有两个可能的纯态”, 而是态空间维数为二. 例如双势阱左右局域制备$L,R$中矩阵为$\left(\begin{smallmatrix}E_L&K\\K^*&E_R\end{smallmatrix}\right)$. 对角能量表左右偏置, 非对角$K$使局域态可改变. 它给$E_c=(E_L+E_R)/2$, $h_z=(E_L-E_R)/2$, $h_x=\operatorname{Re}K,h_y=-\operatorname{Im}K$, 各量都为能量. 厄米$2\times2$矩阵有四个实参数, 可唯一写为$H=E_cI+\boldsymbol h\cdot\boldsymbol\sigma$, $\boldsymbol h\in\R^3$. 由于$(\boldsymbol h\cdot\boldsymbol\sigma)^2=|\boldsymbol h|^2I$, 本征值为$E_c\pm|\boldsymbol h|$. 假定$E_c,\boldsymbol h$不随时间改变. 指数的偶次、奇次项分开求和, 得
\[
U(t)=\e^{-\ii E_ct/\hbar}\left[
\cos\frac{|\boldsymbol h|t}{\hbar}I
-\ii\sin\frac{|\boldsymbol h|t}{\hbar}\frac{\boldsymbol h\cdot\boldsymbol\sigma}{|\boldsymbol h|}\right],
\]
$\boldsymbol h=0$时按连续极限取$U=\e^{-\ii E_ct/\hbar}I$. 公共能量$E_c$只贡献整体相位, 能级差$2|\boldsymbol h|$决定可观测振荡.

\begin{example}
在正交基$\ket L,\ket R$中, $H=E_cI-\Delta\sigma_x$, $\Delta>0$. $\Delta$为能量耦合而非频率, 初态$\ket L$. 求在时间$t$测得$\ket R$的概率.
\end{example}
\begin{solution}
演化给$\ket{\psi(t)}=\e^{-\ii E_ct/\hbar}[\cos(\Delta t/\hbar)\ket L+\ii\sin(\Delta t/\hbar)\ket R]$, 所以$P_R(t)=\sin^2(\Delta t/\hbar)$. 能量本征態为$(\ket L\pm\ket R)/\sqrt2$, 初态同时包含两者. 转移不是沿某条指定经典路径发生; 它由两个本征态的相对相位产生.
\end{solution}

加入不对称项$\epsilon\sigma_z$后, 令$\Omega=\sqrt{\Delta^2+\epsilon^2}/\hbar$, 转移概率为$P_R=\Delta^2\sin^2(\Omega t)/(\Delta^2+\epsilon^2)$. 概率振荡的角频率为$2\Omega$, 周期$\pi/\Omega$, 最大值为$\Delta^2/(\Delta^2+\epsilon^2)$. 能级失谐会阻止完全转移, 这从演化矩阵的非对角元直接得到. 磁共振和氨分子二态模型使用同一代数, 但具体的$\boldsymbol h$须从各自物理相互作用识别.
